\documentclass[12pt,a4paper]{report}

\usepackage{relsize}
\usepackage{graphicx}
\usepackage[margin=3cm]{geometry} 
\usepackage{fancyhdr}
\usepackage{titling} 
\usepackage{mathtools,amsthm,amssymb} 
\usepackage[shortlabels,inline]{enumitem}
\usepackage{indentfirst} 
\usepackage{float}
% \usepackage{setspace}
\usepackage{tikz}\usetikzlibrary{arrows.meta,decorations.pathmorphing}
\usepackage{hyperref}\hypersetup{
    colorlinks=true,
    linkcolor=black, 
    filecolor=blue,
    urlcolor=blue,
    citecolor=cyan
}

\theoremstyle{definition}
\newtheorem{theorem}{Theorem}[section]
\newtheorem{lemma}[theorem]{Lemma}
\newtheorem{corollary}[theorem]{Corollary}
\newtheorem{definition}[theorem]{Definition}
\newtheorem{property}[theorem]{Property}
\newtheorem{proposition}[theorem]{Proposition}
\newtheorem{axiom}[theorem]{Axiom}
\newtheorem{example}[theorem]{Example}
\newtheorem{problem}[theorem]{Problem}

\theoremstyle{remark}
\newtheorem*{remark}{Remark}      
\newtheorem*{solution}{Solution}
\newtheorem*{exercise}{Exercise} 
\newtheorem*{claim}{Claim}

\newcommand{\N}{\mathbb{N}}
\newcommand{\Z}{\mathbb{Z}}
\newcommand{\Q}{\mathbb{Q}}
\newcommand{\R}{\mathbb{R}}
\newcommand{\C}{\mathbb{C}}
\newcommand{\F}{\mathbb{F}}
\newcommand{\rank}{\operatorname{rank}}
\newcommand{\nullity}{\operatorname{nullity}}
\newcommand{\la}{\langle}
\newcommand{\ra}{\rangle}

\renewcommand{\headrulewidth}{0.4pt} 
\renewcommand{\footrulewidth}{0.4pt} 
\setlength{\headheight}{15pt}  
\setlength{\parindent}{22pt}   
\pagestyle{fancy}             
\lhead{\rightmark}              
\chead{}                     
\rhead{}          
\lfoot{\theauthor}        
\cfoot{}
\rfoot{\thepage}              
\linespread{1.5}
% \onehalfspacing
\setcounter{tocdepth}{1} % hide subsections in tableofcontents

\title{Applied Differential Equations}
\author{Ray Chien}
\date{\today}

\begin{document}
\maketitle
\tableofcontents

\chapter{Introductiono to Differential Equations}

\section{Definitions and Terminology}

\begin{definition}[Differential Equations]
    An equation containint the derivatives of one or more unknown functions(or dependent variable), with respect to one or more independent variables, is said to be a \textbf{differential equation(DE).}
\end{definition}

\begin{definition}[Ordinary, Partial]
    If a differential equation contains only ordinary derivatives of one or more unknown function with respect to a single independent variable, it is said to be an \textbf{ordinary differential equation(ODE).} An equation involving partial derivatives of one or more unknown functions of two or more independent variables is called a \textbf{partial differential equation(PDE).}
\end{definition}

\begin{definition}[Order]
    The \textbf{order} of a differential equatin is the order of the highest derivative in the equation.
\end{definition}

\begin{definition}[Linear]
    An ODE is called \textbf{linear} if it has the form 
    \[a_n(x)\frac{d^n y}{dx^n} + a_{n-1}(x)\frac{d^{n-1} y}{dx^{n-1}} + \cdots + a_1(x)\frac{dy}{dx} + a_0(x)y = g(x).\]
\end{definition}

\begin{definition}[Homogeneous Linear ODE]
    A linear ODE is called \textbf{Homogeneous} if it has the form 
    \[a_n(x)\frac{d^n y}{dx^n} + a_{n-1}(x)\frac{d^{n-1} y}{dx^{n-1}} + \cdots + a_1(x)\frac{dy}{dx} + a_0(x)y = 0.\]    
\end{definition}

\begin{remark}
    If $y_1$ and $y_2$ are two solutions of a homogeneous linear ODE, then $y_1 + y_2$ is also a solution of the very ODE.  
\end{remark}

\begin{definition}[$n$-Parameter Family of Solutions, Particular Solution]
    An \textbf{\textit{n}-parameter} family of solutions of an ODE has the form $G(x, y, c_1, c_2, \ldots, c_n) = 0.$ A \textbf{particular solution} is a solution obtained by assigning specific values to these parameters. 
\end{definition}

\begin{definition}[Singular Solution]
    A \textbf{singular solution} of an ODE is a solution that cannot be obtained from the $n$-parameter family of solutions. It only happend in non-linear ODEs.
\end{definition}

\section{Initial-Value Problems}

\begin{definition}[$n$th-Order Initial-Value Problems]
    An \textbf{\textit{n}th-order initial-value problem(IVP)} is an $n$th order ODE and the first $n - 1$ derivatives of $y$ at $x_0.$
\end{definition}

\begin{theorem}[Existence of a Unique Solution]
    For a first-order ODE
    \[\frac{dy}{dx} = f(x, y).\]
    Let $R$ be a rectangular region in the $xy$-plane defined by $a \leq x \leq b, c \leq y \leq d$ that contains the point $(x_0, y_0)$ in its interior. If $f(x, y)$ and $\frac{\partial f}{\partial y}$ are continuous on $R$, then there exists some interval: $I_0(x_0 - h, x_0 + h), h> 0$, contained in $[a, b]$, and a unique function $y(x),$ defined on $I_0$, that is a solutin of the IVP.
\end{theorem}


\chapter{First-Order Differential Equations}

\section{Separable Equations}

\begin{definition}[Saparable]
    A first-order ODE of the form 
    \[\frac{dy}{dx} = g(x)h(y)\]
    is said to be \textbf{separable} or to have separable variables.
\end{definition}

Here are the steps to solve a first order separable ODE:

\begin{enumerate}
    \item We obtain 
    \[\frac{dy}{h(y)} = g(x) \, dy.\]
    \item We let
    \[p(y) = \displaystyle\frac{1}{h(x)}.\]
    \item We do integration on zth sides
    \[\int p(y) \, dy = \int h(x) \, dx.\]
    \item Let $P$ and $H$ be antiderivatives of $p$ and $h,$ we have 
    \[P(y) = H(x) + c,\]
    where $c$ is a constant.
    \item Since the first-order separable ODEs can be non-linear, we may have to find a singular solution. In this case, since we considered $h(x) \neq 0$ in the previous divition, the singular solution occurs when $h(x) = 0.$ Therefore, we check if the solution has the form $y = r$, where $r$ is a constant.
\end{enumerate}

\section{Linear Equations}

\begin{definition}[Standard form of Linear First-Order ODE]
    The \textbf{Standard form} of linear first-ordered differential equaitons is
    \[\frac{dy}{dx} + P(x)y = f(x).\]
\end{definition}

The idea of the method solving first-order differential equations is to multiply the equation with an \textbf{integrating factor} $\mu(x)$ such that the LHS becomes a direct differential of another functon, and the problem thus turns into a simple integraiton.  

\begin{enumerate}
    \item We multiply the standard form of the linear first-order ODE with $\mu(x)$:
    \[\mu(x)\frac{dy}{dx} + \mu(x)P(x)y = \mu(x)f(x).\]
    \item Observe the LHS have the form of product rule, we have
    \[\text{LHS} = \left[\mu(x)y\right]' = \mu(x)\frac{dy}{dx} + \mu'(x)y.\]
    Therefore, we obtain
    \[\mu'(x) = \mu(x) P(x)\]
    \item 
    By solving the equation as a separable equation, we have 
    \begin{align*}
        \ln|\mu(x)| &= \int P(x) \, dx\\
        \mu(x) &= e^{\int P(x) \, dx}
    \end{align*}
    Note that we do not need to consider the integral constant and the sign because they will be eliminated in the next step
    \item We substitude the result into the equation:
    \[[e^{\int P(x) \, dx} y]' = e^{\int P(x) \, dx}f(x) + c\]
    \[y = e^{-\int P(x) \, dx} \left(\int{e^{\int P(x) \, dx}f(x) \, dx} + c\right).\]
    \item Check for singular points.
\end{enumerate}

The other method takes use of the property of the homogeneous linear ODEs. The solution should have the form $y_c + y_p$, where $y_c$ is the family of the solutions to $\displaystyle\frac{dy_c}{dx} + P(x)y_c = 0$, and $y_p$ is a particular solution to $\displaystyle\frac{dy_p}{dx} + P(x)y_p = f(x).$

To proof that the solution of a linear first-order ODE have the form $y = y_c + y_p,$ we write

\[(\frac{dy}{dx} + P(x)y) - (\frac{dy_p}{dx} + P(x)y_p) = \frac{d(y - y_p)}{dx} + P(x)(y - y_p) = 0.\]

\noindent Therefore, $y = y_c + y_p.$

\begin{enumerate}
    \item For $y_c$, we have $\displaystyle\frac{dy_c}{dx} + P(x)y_c = 0$, which is a separable equation having solutions $e^{-\int P(x) \, dx}.$
    \item For $y_p$, we have $\displaystyle\frac{dy_p}{dx} + P(x)y_p = f(x).$ Just notice one or it is actually $e^{-\int P(x) \, dx} \displaystyle\int{e^{\int P(x) \, dx}f(x) \, dx}$.
    \item Check for singlular points
\end{enumerate}

For a given linear first-order ODE, the form may not be standard. Generally, a linear first-order ODE have the form 
\[a_1(x) y' + a_0(x) y = g(x).\]
For $a \neq 0,$ we have
\[y' + \frac{a_0}{a_1} y = \frac{g(x)}{a_1(x)}\]
Therefore, to obtain singular points, solve $a_1(x) = 0.$

\section{Exact Equations}

For exammple, to solve the ODE
\[y \, dx + x \, dy = 0,\]
we notice that the equation is equivelent to
\[d(xy) = 0.\]
Therefore, we obtain the solution $xy = c.$ In this section, we introduce methods to help us to determine whether $M(x, y) \, dx + N(x, y) \, dy$ is the differential of a function $f(x, y)$ and to find it. 

\begin{definition}(Exact Equations)
    A differential expression $M(x, y) \, dx + N(x, y) \, dy$ is an \textbf{exact differential} in a region $R$ of the $xy$-plane if it correspondes to the differential of some function $f(x, y)$ defined in $R$. A first-order differntial of teh form $M(x, y)\,dx + N(x, y)\,dy = 0$ is said to be an \textbf{exact equation} if the expression on the left-hand side is an exact differential.
\end{definition}

\begin{theorem}
    Let $M(x, y)$ and $N(x, y)$ be continuous and have continuous first partial dirivatives in a rectangular region $R$ defined by $a < x < b, c < y < d.$ Then a necesssary and sufficent condition that $M(x, y) \, dx + N(x, y) \, dy$ be an exact differential is 
    \[\frac{\partial M}{\partial y} = \frac{\partial N}{\partial x}.\]
\end{theorem}

\begin{proof}
    It is clear that if the equatino is differential of $f(x, y)$, then
    \[M(x) = \frac{\partial f}{\partial x} \quad \text{and} \quad N(x) = \frac{\partial f}{\partial y}.\]
    If $M, N$ have continuous first partial derivatives, then 
    \[\frac{\partial M}{\partial y} = \frac{\partial}{\partial y}\left(\frac{\partial f}{\partial x}\right) = \frac{\partial N}{\partial x}.\]

    The sufficent part of is a constructive proof, in which we actually find $f$ for $M$ and $N$ that satisfy the condition.
\end{proof}

We now introduce the procedure to find $f$

\begin{enumerate}
    \item By $\displaystyle \frac{\partial f}{\partial x} = M(x, y),$ we have 
    \[f(x, y) = \int M(x, y)\, dx + g(y)\]
    with $g(y)$ indicate the ``constant'' of the integral.
    \item By $\displaystyle \frac{\partial f}{\partial y} = N(x, y)$, we have 
    \[\frac{\partial}{\partial y} \int M(x, y)\, dx +  g'(y) = N(x, y)\]
    \[g(y) = \int \left(-\frac{\partial}{\partial y}\left(\int M(x, y) \, dx\right)  + N(x, y) \right) \, dy.\]
    \item Combining the result of the two steps, we have
    \[f(x, y) = \int M(x, y)\, dx + \int N(x, y) \, dy - \int \left(\frac{\partial}{\partial y}\left(\int M(x, y) \, dx\right)\right) \, dy.\]
\end{enumerate}


The procedure can also be with $x, y$ exchanged.

Since the expression $g'(y)$ can be proved to be independent to be independent to $x$ by
\[\frac{\partial}{\partial x} \left(-\frac{\partial}{\partial y}\left(\int M(x, y) \, dx\right)  + N(x, y) \right) = \frac{\partial M}{\partial x} - \frac{\partial N}{\partial y} = 0.\]
This is also the condition required. Also, this means we can get $f$ by 
\[f(x, y) = \int M_0(x) \, dx + \int N(x, y) \, dx = \int M(x, y) \, dx + \int N_0(y) \, dy\]
where $M_0$ is the part of $M$ which is independent to $y$, and $N_0$ is the part of $N$ wihch is independent to $x$. However, it is not always that easy to seperate the functions. 

Also, $\displaystyle{\int M(x, y) \, dx}$ rigirously means the set of functions whose partial derivatives with respect to $x$ is $M(x, y).$ For example, $xy + h(y) = \displaystyle\int y \, dx$, where $h(y)$ is any function rith respect to $y$ as the ``constant'' of the integral, like $2y$, is correct. Here, however, we separate out the constant part as $g(y)$ in the first step, and ommit the constant as a function with respect to $x$ in the second step.

In addition, some exact function can be obtained by multiplying an integrating factor, that is 
\[\mu(x, y)M(x, y) \, dx + \mu(x, y)N(x, y) \, dy = 0\]
is an exact equation. To find $u$, we check if the condition of the exact equation hold
\[\mu_y M+\mu M_y = \mu_x N + \mu N_x\]
We only look for if $\mu$ is dependent to a single variable, let's say $x$, and independent to the other. We have
\[\mu \frac{M_y - N_x}{N} = \mu_x.\] 
If $\displaystyle\frac{M_y - N_x}{N}$ is independent with $y$, It is relatively easy to solve for it is seperable. The solution is
\[\mu = \int e^\frac{M_y - N_x}{N} \, dx\]
if $\mu$ is dependent to $y$ or 
\[\mu = \int e^\frac{N_x - M_y}{M} \, dy\]
if $\mu$ is dependent to $x$.
Note that if $M_y - N_x = 0$, then $\mu$ is a constant and not needed. Also notice that for linear equations $y' + P(x)y = f(x),$ we have $(P(x)y - F(x))dx - dy = 0.$ We have the integrating factor $\mu = \displaystyle\int e^\frac{M_y - N_x}{N} \, dx = \displaystyle\int e^{P(x)} \, dx.$
In conclustion, apply the method of exact equations if
\begin{enumerate}
    \item $M_y = N_x.$
    \item $\displaystyle\frac{M_y - N_x}{N}$ is independent to $y.$
    \item $\displaystyle\frac{N_y - M_x}{M}$ is independent to $x.$
\end{enumerate}

\section{Solutions By Substitutions}

\begin{definition}[Homogeneous Functions]
    Let $f(x, y)$ be a function. If 
    \[f(tx, ty) = t^\alpha f(x, y)\]
    for any $t \in \R$, then $f$ is said to be \textbf{homogeneous} of degree $\alpha.$
\end{definition}

For a first-order ODE

\[M(x, y) \, dx + N(x, y) \, dy.\]

If $M$ and $N$ are both $\alpha$-degree homogeneous, then let $u = y/x,$ we have 
\[x^\alpha M(1, u) \, dx+ x^\alpha N(1, u) \, dy = 0.\]
With $dy = x\,du + u\,dx$ we have
\[(M(1, u) + uN(1, u)) \, dx + xN(1, u)\,du = 0,\]
which is a seperate equation to solve easily.

\begin{definition}[Bernoulli's Equation]
    If a differential equation have a form
    \[\frac{dy}{dx} + P(x)y = f(x)y^n,\]
    we call it \textbf{Bernoulli's equation}.
\end{definition}
For a Bernoulli's Equation with $n \neq 0$ and $n \neq 1,$ let $u = y^{1 - n}, du = (1 - n)y^{-n}dy$ we have 
\[\frac{y^n}{1 - n}\frac{du}{dx} + P(x)y = f(x)y^n\]
\[\frac{1}{1 - n}\frac{du}{dx} + P(x)u = f(x),\]
which is a linear equation to solve easily.

For equations having the form
\[\frac{dy}{dx} = f(Ax + By + C),\]
with $B \neq 0$, let $u = Ax + By + C,du = A\,dx + B\, dy,$ we have
\[\frac{du}{dx} = -Bf(u) - A,\]
which is a seperated equation to solve easily.

\chapter{Higher-Order Differential Equations}

\section{Linear Equations}

\begin{theorem}
    Let $g(x)$ and $a_i(x)$ for $i = 1, 2, \ldots, n$ be continuous on an interval $I$ including $x_0$ and let $a_n$ for every $x$ in this interval, then the linear ODE
    \[\sum_{i = 0}^{n} a_i(x)\frac{d^ny}{dx^n} = g(x)\]
    has a unique solution if the initial conditions $a_i(x_0)$ are known for $i = 1, 2, \ldots, n.$
\end{theorem}
\begin{proof}
    IDK :(
\end{proof}

Other forms of conditions of ODEs can lead to many, one, or no solutions.  

\begin{definition}[Associated Homogeneous Equations]
    The \textbf{associated homogeneous equation} of a linear ODE
    \[\sum_{i = 0}^{n} a_i(x)\frac{d^ny}{dx^n} = g(x)\]
    is
    \[\sum_{i = 0}^{n} a_i(x)\frac{d^ny}{dx^n} = 0.\]
\end{definition}

\begin{theorem}
    Let $y_1, y_2, \ldots, y_k$ be solutions of a homogeneous differential equation on an intercal $I$, then the linear combination
    \[y = c_1y_1(x) + c_2y_2(x) + \cdots + c_ny_n(x)\]
    where $c_i, i = 1, 2, \ldots, k$ are arbitrary constants, is also a solution on the interval.
\end{theorem}

\begin{proof}
    Trivial.
\end{proof}

\begin{theorem}
    Let $f_1, f_2, \ldots, f_n$ be functions of $x$. If the \textbf{Wronskian}
    \[W(f_1, f_2, \ldots, f_n)
    \begin{vmatrix}
        f_1 & f_2 & \cdots & f_n\\
        f_1' & f_2' & \cdots & f_n'\\
        \vdots & \vdots & \ddots & \vdots\\
        f_1^{(n - 1)} & f_2^{(n - 1)} & \cdots & f_n^{(n - 1)}\\
    \end{vmatrix}
    \]
    is nonzero for all $x$, then the set is linearly independent.
\end{theorem}

\begin{proof}
    IDK :(
\end{proof}

\begin{definition}[Fundamental Set of Solutions]
    A linearly indpendent set of solutions to an $n$th-order ODE on an interval $I$ with cardinality $n$ is said to be a \textbf{fundamental set of solutions} on the interval. 
\end{definition}

\begin{theorem}
    There exists a fundamental set of solutions for any $n$th-order homogeneous linear DE.
\end{theorem}
\begin{proof}
    IDK!!!
\end{proof}

\begin{theorem}
    The general solution of an $n$th-order homogeneous linear equation is the linear combination of the function in the set of its fundamental set of solutions.
\end{theorem}
\begin{proof}
    IDKIDKIDK
\end{proof}

\section{Reduction of Order}

If $y_1(x)$ is found as a solution of a second-order ODE 
\[a_2(x)y'' + a_1(x)y' + a_0(x)y = 0,\]
we can reduce order by letting $y_2(x) = u(x)y_1(x).$ First, we transform the equation into standard form
\[y'' + P(x)y' + Q(x)y = 0.\]
With $y = uy_1, y' = u'y_1 + uy_1', y_1'' = u''y_1 + 2u'y_1' + uy_1''$, we have
\[u(y_1'' + Py_1' + Qy_1) + u''y_1 + u'(2y_1' + Py_1) = u''y_1 + u'(2y_1' + Py_1) = 0.\]
Let $w = u',$ we have 
\[w'y_1 + w(2y_1' + Py_1) = 0,\]
\[\frac{dw}{w} + 2\frac{y_1'}{y_1}\, dx + P \, dx\]
\[\frac{dw}{w} + 2\frac{dy_1}{y_1} + P \, dx,\]
from which we obtain
\[wy^2 = c_1e^{\int - P \, dx}\]
\[u = c_1\int\frac{e^{\int -P \, dx}}{y_1^2} \, dx + c_2.\]
By choosing $c_1 = 1$ and $c_2 = 0$, we have 
\[y_2 = y_1 \int \frac{e^{\int -P \, dx}}{y_1^2} \, dx.\]

\section{Homogeneous Linear Equations with Constan}

\begin{definition}[Auxiliary Equation]
    For a homogeneous linear equation with constant coefficients
    \[\sum_{i = 0}^{n} a_iy^{(i)} = 0,\]
    its corresponding \textbf{auxiliary equation} is
    \[\sum_{i = 0}^{n} a_i m^i = 0.\]
\end{definition}

\begin{remark}
    The auxiliary equation is obtained by letting $y = e^{mx}.$
\end{remark}

\begin{theorem}
    For a homogeneous linear equation with constant coefficients, let $m_1, m_2,\ldots, m_k$ be $k$ distinct solution of its auxiliary equation, and $t_i$ be the multiplicity of $m_i,$ then the general solution of the homogeneous linear equation is
    \[\sum_{i = 1}^{k} \sum_{j = 0}^{t_i - 1} c_{ij}x^{t_j} e^{m_i}x\]
    where $c_{ij}$ are constants.
\end{theorem}

\begin{proof}
    Just looks damn right.
\end{proof}

\begin{remark}
    It is often to see $c_1e^{\alpha + i\beta} + c_2e^{\alpha - i\beta}$ included in the expresstion of the solution, where $i = \sqrt{-1}.$ We can obtain cleaner expresstion with $c_1e^{\alpha + i\beta} + c_2e^{\alpha - i\beta} = e^\alpha(c_1 (\cos \beta + i \sin \beta) + c_2 (\cos \beta + i \sin \beta)) = e^\alpha(a \cos \beta + b \cos \beta),$ where $a$ and $b$ are constant, might be complex numbers.
\end{remark}

\section{Undetermined Coefficients -- Superposition Approach}

This method is less or more by experience. The approach can deal with nonhomogeneous linear DEs with constant coefficients, that is
\[\sum_{i = 0}^{n} a_iy^{(i)} = g(x)\]
where $g + g' + g'' +\ldots$ must have finite terms. We first devide $g(x)$ to different kinds of functions, such as polynomial, exponential, a sine or cosine function, or products of them. The particular solution of the DE is the superposition of the functions with undetermined coefficients.   

Sometimes, the guessed particular solution may be within the homogeneous solution. Whenever that happens, just times $x$.

\begin{table}[H]
\centering
\caption{Trial Particular Solutions}
\label{tab:trial_solutions}
\begin{tabular}{ll}
\hline
$g(x)$ & Form of $y_p$ \\
\hline
$67$ & $A$ \\
$5x + 7$ & $Ax + B$ \\
$3x^2 - 2$ & $Ax^2 + Bx + C$ \\
$x^3 - x + 1$ & $Ax^3 + Bx^2 + Cx + E$ \\
$\sin 4x$ & $A\cos 4x + B\sin 4x$ \\
$\cos 4x$ & $A\cos 4x + B\sin 4x$ \\
$e^{5x}$ & $Ae^{5x}$ \\
$(9x - 2)e^{5x}$ & $(Ax + B)e^{5x}$ \\
$x^2 e^{5x}$ & $(Ax^2 + Bx + C)e^{5x}$ \\
$e^{3x}\sin 4x$ & $Ae^{3x}\cos 4x + Be^{3x}\sin 4x$ \\
$5x^2\sin 4x$ & $(Ax^2 + Bx + C)\cos 4x + (Ex^2 + Fx + G)\sin 4x$ \\
$xe^{3x}\cos 4x$ & $(Ax + B)e^{3x}\cos 4x + (Cx + E)e^{3x}\sin 4x$ \\
\hline
\end{tabular}
\end{table}


\section{Undetermined Coefficients -- Annihilator Approach}

This method requires the same condition to the last one.

\begin{definition}[Annihilator]
    If $L$ is a linear differential operator with constant coefficients and $f$ is a sufficiently differentiable function such that 
    \[L(f(x)) = 0,\]
    then $L$ is said to be an \textbf{Annihilator} of $f$.
\end{definition}
\begin{remark}
    For constant coefficient linear operators $L_1, L_2, L_3,$ we have 
    \begin{enumerate}
        \item $(L_1L_2)y = L_1(L_2y)$
        \item $L_1L_2 = L_2L_1$
        \item $(L_1L_2)L_3 = L_1(L_2L_3).$
    \end{enumerate}
    Also, Let $L_4$ and $L_5$ be annihilators of $f$ and $g$, then $L_4L_5(f + g) = 0.$
\end{remark}

Here are some general solutions to find annihilator.

\begin{enumerate}
    \item $D^n$ for $n - 1$-ordered polynomials.
    \item $(D - \alpha)^n$ for $\displaystyle\sum_{i = 0}^{n - 1} c_i x^i e^{\alpha x}.$
    \item $(D^2 + 2\alpha D + \alpha^2 + \beta^2)^n$ for $\displaystyle\sum_{i = 0}^{n - 1} (a_i x^i e^{\alpha x} \cos \beta + b_i x^i e^{\alpha x} \sin \beta)$.
\end{enumerate}

Now, for a nonhomogeneous linear DE with constant coefficients, that is
\[L(y)= g(x),\]
with $g + g' + g'' +\ldots$ having finite terms, and $L_1$ is an annihiliator of $g$, then the particular solution is the solution with undetermined coefficient of the DE is the homogeneous solution of $L_1Ly.$

\section{Variation of Parameters}

In this section we focus on how to get particular solutoins of linear ODEs by given homogeneous solutions, in general cases. Let us take a look at a linear equation again. Consider
\[y^{(n)} + p_{n - 1}(x)y^{(n - 1)} + \cdots + p_0(x)y = f(x).\]
Let
\[\mathbf{Y_k} = \begin{pmatrix}
    y_k\\ y_k' \\ \vdots \\ y_k^{(n - 1)}
\end{pmatrix}, \quad \mathbf{\Phi} =  \begin{pmatrix}
    \mathbf{Y_1} & \mathbf{Y_2} & \cdots & \mathbf{Y_n} 
\end{pmatrix}\]
where $k = 1, 2, \ldots, n$ and $y_1, y_2, \ldots, y_n$ are linearly indpendent solutions of the corresponding homogeneous equation, that is
\[
\mathbf{Y_k} = \mathbf{A}\mathbf{Y_k}, \quad \text{so} \quad 
\mathbf{\Phi}' = \mathbf{A}\mathbf{\Phi}
\]
where
\[\mathbf{A}=\begin{pmatrix}0&1&0&\cdots&0\\0&0&1&\cdots&0\\\vdots&&&\ddots&\vdots\\0&0&0&\cdots&1\\-p_0&-p_1&-p_2&\cdots&-p_{n-1}\end{pmatrix}.\]
Now, consider a particular solution
\[\mathbf{Y_p} = \mathbf{A}\mathbf{Y_p} + \mathbf{F},\]
where
\[\mathbf{F} = \begin{pmatrix}
    0 \\ 0 \\ \vdots \\ f(x) 
\end{pmatrix}.\]
Now, since $\mathbf{\Phi}$ is invertable, let
\[\mathbf{u} = \begin{pmatrix}
    u_1(x) & u_2(x) & \cdots & u_n(x)
\end{pmatrix}= \mathbf{\Phi}^{-1}\mathbf{Y_p},\]
so
\[
\mathbf{Y_p} = \mathbf{\Phi}\mathbf{u},
\]
that is, $y_p = u_k(x)y_k(x)$ for $k = 0, 1, \ldots n$. Next, we write the ODE again
\[\mathbf{Y_p}' = \mathbf{A}\mathbf{Y_p} + \mathbf{F}\]
\[\mathbf{\Phi}'\mathbf{u} + \mathbf{\Phi}\mathbf{u}' = \mathbf{A} \mathbf{\Phi}\mathbf{u} + \mathbf{F}\]
\[\mathbf{A}\mathbf{\Phi}\mathbf{u} + \mathbf{\Phi}\mathbf{u}' = \mathbf{A} \mathbf{\Phi}\mathbf{u} + \mathbf{F}\]
\[\mathbf{\Phi}\mathbf{u}' =  \mathbf{F}.\]
Therefore, we can obtain $u_k$ by Cramer's rule and find a particular solution.

\section{Cauchy-Euler Equations}
\begin{definition}[Cauchy-Euler Equation]
    A linear ODE of the form 
    \[\sum_{k = 1}^{n} a_k x^k y^{(k)} = g(x)\]
    where $a_k$ are constant coefficients, is called a \textbf{Cauchy-Euler equation}.
\end{definition}

For Cauchy-Euler equations, let $y = x^m$ and solve the auxiliry equation about $m$. For repeating root of, that's say $m_1$ with multiplicity $l$, the corresponding $l$ independent solutions are $x^{m_1}, x^{m_1}\ln x, \ldots, x^{m_1}(\ln x)^l$. For complex roots, $\alpha \pm i\beta$, the corresponding two solutions are $x^\alpha \cos(\beta \ln x)$ and $x^\alpha \sin(\beta \ln x)$.

Notice that it seems to be similar to the rules of constant linear ODE with coeffecients but now $x$ becomes $\ln x$. That is because the whole method for Cauchy-Euler equations mentioned can also be done by substitute $x = e^t$ and get exact same auxiliry functions.

The substitution of $x = e^t$ also help us get particular solutions of Cauchy-Euler equations. For example, consider
\[x^2 y'' - xy' + y = \ln x.\]


\chapter{Series Solutions of Linear Equations}

\section{Solutions About Ordinay Points}

The idea is that, consider a linear ODE
\[\sum_{i = 0}^{n} a_i(x)y^{(i)}(x) = 0.\]
If the solution is analytic, that is, can be expressed as
\[y_i = \sum_{k = 0}^{\infty} c_k (x - a)^k,\]
then we can obtain the coeffiecients by $a_i(x)$ if they are also analytic. If $a_n = 0,$ then since $x - a$ is a factor of $a_n(x)$, we actually only have $n - 1$th order recursive equation of $c_k$, that is, cannot produce $n$ linearly independent solutions.

\begin{definition}[Analytic]
    Let $f$ be a function. If we can write
    \[f(x) = \sum_{k = 0}^{\infty} a_k(x - x_0)^k\]
    on some open interval including $x_0$, then $f$ is \textbf{analytic} at $x_0$.
\end{definition}

\begin{definition}[Ordinary Points and Singular Points]
    For a linear ODE
    \[\sum_{i = 0}^{n} a_i(x)y^{(i)}(x) = 0.\]
    Let
    \[P_k(x) = \frac{a_k(x)}{a_n(x)}\]
    for $k = 0, 1, 2, \ldots, n - 1.$ Then $x_0$ is an \textbf{ordinary point} if $P_0, P_1, \ldots, P_{n - 1}$ are analytic at $x_0$. Otherwise, it is a \textbf{singular point.}
\end{definition}

\begin{theorem}
    Consider 
    \[\sum_{i = 0}^{n} a_i(x)y^{(i)}(x) = 0,\]
    where
    $a_0, a_1, \ldots, a_n$ be functions of $x$ analytic at $a$ and $a_n(a) \neq 0$. 
    There exist $n$ linearly independent analytic solutions around $a$ converge at least for $|x - a| < R,$ where $R$ is the distance from $a$ to the nearst singular point.
\end{theorem}

Before jump into the next chapter to examine a particular method for singular points, let us try three methods to show that singular points do matters.

\begin{example}
    Consider the ODE
    \[y' - x^2y = 0.\]
    Let
    \[x = \sum_{k = 0}^{\infty} c_k x^{k}\]
\end{example}

\begin{example}
    \[y' - y = 0\]
\end{example}



\begin{example}
    \[xy' + y = 0\]
    where $x = 0$ is a singular point. Let us whatever do this to see the problem.
    \[x \sum_{k = 1}^{\infty}kc_kx^{k - 1} + \sum_{k = 0}^{\infty} c_k x^{k} = 0\]
    \[\sum_{k = 0}^{\infty}{(k + 1)}c_{k}x^{k} + \sum_{k = 0}^{\infty} c_k x^{k} = 0.\]
    This gives us $y = 0$ which is complete useless. However, if we consideer to shift the power of $x$ by $r$, that is
    \[x = \sum_{k = 0}^{\infty} c_k x^{k + r}\]
    we have
    \[x \sum_{k = 0}^{\infty}{(r + k)}c_kx^{r + k - 1} + \sum_{k = 0}^{\infty} c_k x^{r + k} = 0\]
    \[\sum_{k = 0}^{\infty}(r + k)c_{k}x^{r + k} + \sum_{k = 0}^{\infty} c_k x^{r + k} = 0\]
    \[(r + k + 1)c_k = 0, k \geq 0.\]
    Notice that this solution reduced to $cx^{-1}.$ For example, let $c_0 \neq 0,$ and others are $0$, then $r = -1$. We have $y = c_0 x^{-1}$. Or, let $c_{67} \neq 0$ and others are $0$, then $r = -68.$ We have $y = c_{67}x^{-68 + 67} = c_{67} x^{-1}.$
\end{example}

\begin{example}
    Consider the ODE
    \[x^2y' + y = 0\]
    where $x = 0$ is a singular point. If we take lessons from the last example, that is
    \[x = \sum_{k = 0}^{\infty} c_k x^{k + r}\]
    we have
    \[x^2 \sum_{k = 0}^{\infty}{(r + k)}c_kx^{r + k - 1} + \sum_{k = 0}^{\infty} c_k x^{r + k} = 0\]
    \[\sum_{k = 1}^{\infty}(r + k - 1)c_{k - 1}x^{r + k} + \sum_{k = 0}^{\infty} c_k x^{r + k} = 0\]
    \[
    \begin{cases}
        (r + k + 1)c_{k + 1} + c_k = 0, k \geq 0\\
        c_0 = 0, k = 0.
    \end{cases}
    \]
\end{example}


\end{document}